\chapter{Elementare Funktionen I}

\setcounter{section}{7}
\setcounter{subsection}{0}
\setcounter{subsubsection}{1}
\setcounter{secnumdepth}{4}
\subsection{Polynome}
\href{https://youtu.be/b0zNTRgnXPM}{Video 01}\\
Ein Polynom $ p: \mathbb{C}\longmapsto  \mathbb{C} $ (oder $\mathbb{R}\longmapsto \mathbb{R}  $) mit:\\
\begin{align*}
p(x)=a_0+a_1\cdot x+a_2\cdot x^2+a_3\cdot x^3+\ldots+a_n\cdot x^n
\end{align*}
wobei:
\begin{quote}
	\textbullet\;\;$ a_0,\ldots,a_n \in \mathbb{C}$ (oder $ \mathbb{R} $) heißen Koeffizienten \\
	\textbullet\;\;$ a_n\neq 0 $ heißt $ a_n $ der höchste Koeffizient\\
	\textbullet\;\;Ist $ a_n\neq 0 $ dann hat $ p $ den Grad $ n $, man schreibt: $ \deg(p)=n\in\mathbb{N} $\\
	\textbullet\;\;Für $ a_0=a_1=\ldots=a_n=0 $ ($ \Leftrightarrow p(x)=0\;\forall x\in \mathbb{C}$) definieren wir: $ \deg(P)=-\infty $\\
	\textbullet \;\;$ x_i\in  \mathbb{C} $ heißt Nullstelle von $ p $, wenn $ p(x_i)=0 $ dann gilt: $ (x-x_i) $ teilt $ p(x) $
\end{quote} 
zum Beispiel:\\
\begin{tabular}[t]{ll}% l=linksbündig r=rechtsbündig
	$\blacktriangleright$ & $ \deg(7+\mathrm{i})=\deg((7+\mathrm{i})\cdot x^0)=0 $  \\
	& Dies ist eine konstante Funktion und auch ein Polynom mit $ x^0 $\vspace*{2mm} \\
	$ \blacktriangleright $	& $ p(x)=7x^2+2 $\\
	&Höchste Koeffizient ist 7\\
	&Der Grad von $ p=\deg(2) $\vspace*{2mm}\\
	$\blacktriangleright$ & $\overset{\AC}{p}(x)=x^2-3x+2$\\
	&Es gilt  $\overset{\AC}{p}(1)=1^2-3\cdot 1+2=0$\\
	&$\Rightarrow x_1=1$ ist Nullstelle von $\overset{\AC}{p}(x)  $\\
	&$ (x-x_1)=(x-1)$ teilt das Polynom $\overset{\AC}{p}(x)  $\\
	&Und es gilt  $\overset{\AC}{p}(2)=2^2-3\cdot 2+2=0$\\
    &$\Rightarrow x_2=2$ ist Nullstelle von $\overset{\AC}{p}(x)  $\\
  &$ (x-x_2)=(x-2)$ teilt das Polynom $\overset{\AC}{p}(x)  $\\
& Warum teilt $ (x-x_i) $ das Polynom $\overset{\AC}{p}(x)  $?\\
&Es gilt: $ \overset{\AC}{p}(x)=(x-1)(x-2)=x^2+3x+2 $\\
&d.h. $ \frac{\overset{\AC}{p}(x)}{(x-1)}=\frac{(x-1)(x-2)}{(x-1)}=(x-2) $
\end{tabular} \\
\subsubsection{Beispiel zum Thema Polynome}
\href{https://youtu.be/cbCiEF4aO6E}{Video 02}\\
Mit Polynome kann man wie mit ganzen Zahlen  rechnen.\\
\newpage
\noindent
Beispiele:\\
\textbullet\;\;Sei $ p_1(x) =2x^2+7\mathrm{i}x+5$ und $\;\; p_2(x)=4x^{12}+7x$\\
\begin{quote}


\begin{align*}
p_1(x)+p_2(x)=(2x^2+7\mathrm{i}x+5)+(4x^{12}+7x)&=4x^{12}+2x^2+7\mathrm{i}x+7x+5\\&=4x^{12}+2x^2+(7\mathrm{i}+7)x+5\\
\end{align*}
Bemerkung: \\
\begin{align*}
&\deg(p_1(x))=2\\
&\deg(p_2(x))=12\\
&\deg(p_1(x)+p_2(x))=12
\end{align*}
\end{quote}
\textbullet \;\; Sei $ p(x)=(7x+1) $ und $ q(x)=(2x^2+2x) $\\
\begin{quote}
	\begin{align*}
	(7x+1)(2x^2+2x)=14x^3+14x^2+2x^2+2x=14x^3+16x^2+2x\\
	\end{align*}
	Bemerkung:\\
	\begin{align*}
&	\deg(p(x))=1\\
&\deg(q(x))=2\\
&\deg(q(x)\cdot p(x))=3\\
&\text{Der Grad des Produktes ist die Summe der Grade der Polynome $ p(x), q(x) $}
	\end{align*}
\end{quote}
\textbullet\;\; Sei $ p(x)=\frac{(x^2+1)}{(x-\mathrm{i})} $\\
\begin{align*}
\frac{(x^2+1)}{(x-\mathrm{i})}=\frac{(x-\mathrm{i})(x+\mathrm{i})}{(x-\mathrm{i})}=x+1
\end{align*}
 
 Bemerkung: $ (x-\mathrm{i})(x+\mathrm{i})=x^2+\mathrm{i}x-\mathrm{i}x-\mathrm{i}^2=x^2+1 $\\
\subsection{Faktorisierung von Polynomen}
\subsubsection{Komplexe Polynome}
\href{https://youtu.be/hK5Ggr4f8As}{Video 03}\\
\begin{center}
	

\fbox{\parbox  {95mm} {
	\begin{center}
		Fundamentalsatz der Algebra\\
		\begin{align*}
		& \text{Sei }\\
		&p: \mathbb{C}\longmapsto  \mathbb{C}\;\;\text{,mit }\;\;p(x)=a_0+a_1 x+\ldots +a_nx^n \\
		&\text{ein Polynom, dann gibt es: }\\
		&x_1,\ldots,x_n\in  \mathbb{C}\;\;\text{s.d. }\\
		&p(x)=a_n(x-x_1)(x-x_2)\cdotp\ldots\cdotp (x-x_n)\;\;\text{wenn }\;\;\deg(p)\geq 1
		\end{align*}
	\end{center}
	}}

\end{center}
Bemerkung:\\
Dies wird Faktorisierung von $ p $ genannt.\\
$ (x-x_n) $ nennt man Linearfaktoren.\vspace*{2mm}\\
Offensichtlich gilt:\\
\begin{align*}
p(x_i)=0\;\;\text{für }\;\;i=1,\ldots,n
\end{align*}
Also sind $ x_i $ mit $ i\in\{1,\ldots,n\} $\;\;Nullstellen von $ p(x) $\vspace*{2mm}\\
Bemerkung:\\
Nullstellen müssen nicht verschieden sein. (z.B: doppelte Nullstelle).\vspace*{2mm}\\
Beispiele:\\
\textbullet\;\;Sei $ p(x)=x^2+1$
\begin{align*}
&x^2+1=(x-\mathrm{i})(x+\mathrm{i})\\
&\text{Nullstellen sind:}\;\;x_1=\mathrm{i}, x_2=-\mathrm{i}\\
\end{align*}
\textbullet\;\;Sei $p(x)=(x^2+1)^2$\\
\begin{align*}
(x^2+1)^2=(x^2+1)(x^2+1)&=(x-\mathrm{i})(x+\mathrm{i})(x-\mathrm{i})(x+\mathrm{i})\\
&=(x-\mathrm{i})^2(x+\mathrm{i})^2
\end{align*}

\subsubsection{Reelle Polynome}
\href{https://youtu.be/yTjK64ypuI0}{Video 04}\\
Sei $ p:\mathbb{R}\longmapsto \mathbb{R} $ mit $ p(x)=a_0+a_1 x+\ldots +a_n x $ ein Polynom mit $ \deg(p)\geq 1 $.\\
Dann existieren:\\
$ x_1,\ldots.x_n\in\mathbb{R} $ mit $ k\leq n $ und ein Polynom $ \overset{\AC}{p}(x) $ mit:\\
\begin{align*}
p(x)=\overset{\AC}{p}(x)\cdot(x-x_1)(x-x_k)
\end{align*}
$\blacktriangleright\;\; x_1,\ldots,x_n $ sind reelle Nullstellen von p.\\
$ \blacktriangleright $\;\; Es kann sein, dass ein Polynom übrig bleibt.\\
$ \blacktriangleright $\;\; Dabei gilt: $ \deg( \overset{\AC}{p}(x)) =n-k$\vspace*{4mm}\\ 
Beispiele:\\
\textbullet \;\;Sei $p:\mathbb{R}\longmapsto\mathbb{R}$ mit $ p(x)=x^2+1 $\\
\begin{align*}
 p(x)=x^2+1=x^2+1\\
\end{align*}
Bemerkung:\\
$ p(x)=x^2+1=(x-\mathrm{i})(x+\mathrm{i}) $ hat zwei komplexe Nullstellen!\\
Deswegen ist dies die faktorisierte Form. Das Polynom $ p(x) $ kann nicht weiter faktorisiert werden.\vspace*{4mm}\\
\textbullet\;\; Sei $p:\mathbb{R}\longmapsto\mathbb{R}$ mit $ p(x) =x^4-1$\\
\begin{align*}
p(x) =x^4-1=(x^2-1)(x^2+1)=(x-1)(x+1)(x^2+1)\\
\end{align*}
Faktorisieren heißt:\\
Mit reellen Polynomen die reelen Nullstellen und mit komplexen Polynomen alle Nullstellen bestimmen.\\
\subsubsection{Nullstellenberechnung von Polynomen}
\paragraph{Nullstellenberechnung von Polynomen vom Grad 2}$  $\vspace*{3mm}\\
\href{https://youtu.be/l27JVKX78Kc}{Video 05}\\
Betrachte:\;\;$ p(x)=2x^2-6\mathrm{i}x-4 $\\
Gesucht: $ x_1 $ und $x_2$ mit $ p(x_1)=p(x_2)=0 $\vspace*{2mm}\\
Mit $ p(x)=0 $ erhalten wir:\\
\begin{align*}
2x^2-6\mathrm{i}x-4&=0\\
2(x^2-3\mathrm{i}x-2)&=0
\end{align*}
Mit der $ p-q- $Formel gilt:\\
\begin{align*}
x_{1,2}=-\frac{p}{q}\pm\sqrt{\bigg(\frac{p}{2}\bigg)^2-q}
\end{align*}
\begin{align*}
x_{1,2}&=-\frac{-3\mathrm{i}}{2}\pm\sqrt{\bigg(\frac{3\mathrm{i}}{2}\bigg)^2-(-2)}=-\frac{-3\mathrm{i}}{2}\pm\sqrt{\bigg(\frac{3\mathrm{i}}{2}\bigg)^2+2}\\
&=+\frac{3\mathrm{i}}{2}\pm\sqrt{\frac{-9+8}{4}}=\frac{3\mathrm{i}}{2}\pm\sqrt{-\frac{1}{4}}\\
\end{align*}
Bemerkung:\\
$ x_{1,2} \in  \mathbb{C}\backslash \mathbb{R}\Rightarrow$ Es existiert keine reelle Nullstelle!\\
$ i^2=-1\Leftrightarrow i=\sqrt{-1} $\\
\begin{align*}
x_{1,2}&=\frac{3\mathrm{i}}{2}\pm\sqrt{-\frac{1}{4}}=\frac{3\mathrm{i}}{2}\pm\sqrt{\frac{1}{4}\cdot(-1)}=\frac{3\mathrm{i}}{2}\pm\sqrt{\genfrac{}{}{}{2}{1}{4 }}\sqrt{-1}\\
&=\frac{3\mathrm{i}}{2}\pm\frac{1}{2}\mathrm{i}\\
\end{align*}
\begin{align*}
&x_1=\frac{3}{2}\mathrm{i}+\frac{1}{2}\mathrm{i}=2\mathrm{i}\\
&x_2=\frac{3}{2}\mathrm{i}-\frac{1}{2}\mathrm{i}=\mathrm{i}\\
\end{align*}
Oder mit quatratischer Ergänzung:\\
\begin{align*}
x^2-3\mathrm{i}x-2&=0\\
x^2-3\mathrm{i}x+\bigg(\frac{-3\mathrm{i}}{2}\bigg)^2-\bigg(\frac{-3\mathrm{i}}{2}\bigg)^2-2&=0\\
x^2-3\mathrm{i}x+\frac{9\mathrm{i}^2}{4}-\frac{9\mathrm{i}^2}{4}-2&=0\\
\bigg(x-\frac{3}{2}\mathrm{i}\bigg)^2&=-\frac{9}{4}+2\\
x-\frac{3}{2}\mathrm{i}&=\pm\sqrt{-\frac{1}{4}}\\
x_{1,2}&=\frac{3}{2}i\pm\sqrt{-\frac{1}{4}}\\
x_1&=2\mathrm{i}\\
x_2&=\mathrm{i}\\
\end{align*}
Fazit:\\
\begin{align*}
p(x)=2(x^2-3\mathrm{i}x-2)=2(x-x_0)(x-x_1)=2(x-2\mathrm{i})(x-\mathrm{i}) 
\end{align*}
\paragraph{Nullstellenberechnung von Polynomen vom Grad größer 2}$  $\vspace*{3mm}\\
\href{https://youtu.be/Phr9F9GFmws}{Video 06}\\
Sei ein Polynom $ p $  vom Grad $ n>2 $ gegeben.\\Gesucht sind die Nullstellen.\vspace*{2mm}\\
Vorgehensweise:\\
Bestimme mittels raten oder ausprobieren alle Nullstellen $ n-2 $ ($x_1, x_2,\ldots,x_{n-2}$).\\
Dann gilt:\\
\begin{align*}
p(x)=\overset{\AC}{p}(x)\cdot (x-x_1)\cdot\ldots\cdot(x-x_{n-2})
\end{align*}
Wegen $ \deg(p)=n $ muss $ \deg(\overset{\AC}{p})=n-(n-2)=2 $ sein.\\
Dann können wir auf $\overset{\AC}{p}  $ die p-q-Formel anwenden.\\
Dabei ist:\\
\begin{align*}
\overset{\AC}{p}(x)=\frac{p(x)}{(x-x_1)\cdot\ldots\cdot(x-x_{n-2})}
\end{align*}
Beispiel:\\
Gegeben das Polynom:   \\
\begin{align*}
p(x)=2x^3-4\mathrm{i}x^2+2x-4\mathrm{i}
\end{align*}
Bestimme die Nullstellen.\\
1. Rate eine Nst.:
\begin{quote}
	$x_1=-\mathrm{i}$\\
	$ x_1$ einsetzen \\
	\begin{align*}
	&2(-\mathrm{i})^3-4\mathrm{i}(-\mathrm{i})^2+2(-\mathrm{i})-4\mathrm{i}=2\mathrm{i}+4\mathrm{i}-2\mathrm{i}-4\mathrm{i}=0\\
	&\Rightarrow \frac{p(x)}{(x+\mathrm{i})}
	\end{align*}
\end{quote}
2. Polynomdivision:\\
\begin{quote}
	\begin{align*}
	& (2x^3-4\mathrm{i}x^2+2x-4\mathrm{i})\div(x+\mathrm{i})=2x^2-6\mathrm{i}x-4\\
	-& \underline{(2x^3+2\mathrm{i}x^2)}\\
	& \hspace{0.7cm}-6ix^2+2x\\
	&\hspace{0.2cm}-\underline{(-6\mathrm{i}x^2+6x)}\\
	&\hspace{1.9cm}-4x-4\mathrm{i}\\
	&\hspace{1.4cm}-\underline{(-4x-4\mathrm{i})}\\
	&\hspace{3.5cm} 0
	\end{align*}
	D.h.:\\
	\begin{align*}
	p(x)&=(2x^2-6\mathrm{i}x-4)\cdot(x+\mathrm{i}) \;\; \text{\tiny mit p-q-Formel \normalsize }\\
	&=2(x-\mathrm{i})(x-2\mathrm{i})(x+\mathrm{i})
	\end{align*}
	Die Nullstellen sind:\\
	$ x_1=-\mathrm{i} $\\
	$ x_2=\mathrm{i} $\\
	$ x_3=2\mathrm{i} $\\
	
	
	
\end{quote}
\subsubsection{Komplexe Nullstellen reeller Polynome}
\href{https://youtu.be/uVsYwAYrcUs}{Video 07}\\
Wenn ein Polynom mit reellen Koeffizienten eine komplexe Nullstelle $ x_1 $ hat, dann ist das komplex konjungierte von $ x_1 $ auch eine Nullstelle.\\
Beispiel:\\
\begin{align*}
p(x)=x^2+1=(x-\mathrm{i})(x+1)\\
\end{align*}
Also Nullstellen sind:\\
\begin{align*}
 x_1&=\mathrm{i} \\
 x_2&=-\mathrm{i}=\overline{x}_1=\overline{\mathrm{i}} \\
\end{align*}
Bemerkung:\\
Komplex konjungiert heißt:\\
\begin{align*}
\overline{a+\mathrm{i}b}&=a-ib\\
\overline{i}&=\overline{0+1\cdot \mathrm{i}}=0-\mathrm{i}=-\mathrm{i}
\end{align*}
Allgemein gilt:\\
Sei $ p: \mathbb{C}\longmapsto  \mathbb{C} $ mit $ p(x)=a_0+a_1x+\ldots +a_nx^n $ mit $ a_0, a_1,\ldots, a_n\in \mathbb{R} $ und ist $ x_1 $ eine komplexe Nullstelle von $ p $, dann ist $x_2= \overline{x}_1 $ auch eine Nullstelle von $ p $.\vspace*{2mm}\\
Beispiel:\\
\begin{align*}
p(x)=x^2+4&=(x-2\mathrm{i})(x+2\mathrm{i})\\
x_1&=2\mathrm{i}\\
x_2&=\overline{x}_1=-2\mathrm{i}\\
\end{align*}
\paragraph{Aufgabe zur Nullstellenberechnung Teil I}$  $\vspace*{3mm}\\
\href{https://youtu.be/ujBUQtLfSI4}{Video 08}\\
Faktorisiere das Polynom $ P(x) $:\\
\begin{align*}
&p(x)=x^7-4x^6+8x^5-11x^4+11x^3-8x^2+4x-1\\
&\text{mit Nst.: }\;\;x_1=1\;\;\text{und }\;\;x_2=\mathrm{i}\\
\end{align*}
Da die Koeffizienten von $ P(x) $ reell sind, ist $ x_3=\overline{x}_2=-\mathrm{i} $ auch eine Nullstelle.\\
Also teilt $ P(x) $:\\
\begin{align*}
(x-1)(x-\mathrm{i})(x+\mathrm{i})=(x-1)(x^2+1)=x^3-x^2+x-1\\
\end{align*}
\begin{align*}
&(x^7-4x^6+8x^5-11x^4+11x^3-8x^2+4x-1)\div(x^3-x^2+x-1)=x^4-3x^3+4x^2-3x+1\\
-&\underline{(x^7-\;\;x^6+\;\;x^5-\;\;x^4)}\\
&\hspace{0.35cm}0-\;3x^6+7x^5-10x^4+11x^3\\
&\hspace{0.22cm}-(\underline{-3x^6+3x^5-\;\;3x^4+\;3x^3})\\
&\hspace{1.62cm} 0+4x^5-\;\;7x^4+\;8x^3-8x^2\\
&\hspace{1.35cm}-\underline{(-4x^5-\;\;4x^4+\;4x^3-4x^2)}\\
&\hspace{2.8cm} 0-\;\;3x^4+\;4x^3-4x^2+4x\\
&\hspace{2.75cm}-(\underline{-3x^4+\;3x^3-3x^2+3x})\\
&\hspace{4.2cm} 0+\;\;x^3-\;\;x^2+\;\;x-1\\
&\hspace{3.85cm}-\underline{(+\;\;\;x^3-\;\;x^2+\;\;x-1)}\\
&\hspace{8.2cm}0
\end{align*}
Dann gilt:\\
\begin{align*}
\frac{P(x)}{(x^3-x^2+x-1)}=x^4-3x^3+4x^2-3x+1
\end{align*}
\begin{align*}
\Leftrightarrow P(x)=(x^4-3x^3+4x^2-3x+1)(x-1)(x-i)(x+i)\\
\end{align*}
\paragraph{Aufgabe zur Nullstellenberechnung Teil II}$  $\vspace*{3mm}\\
\href{https://youtu.be/iHXUOPQJfgM}{Video 09}\\
Faktorisiere $ P_2(x)=x^4-3x^3+4x^2-3x+1$\\
Suche Nullstelle:\\
\begin{align*}
P_2(1)=(1-3+4-3+1)=0
\end{align*}
Also ist $ (x-1) $ teiler von $ P_2(x) $\\
\begin{align*}
&(x^4-3x^3+4x^2-3x+1)\div(x-1)=x^3-2x^2+2x-1\\
-&\underline{(x^4-x^3)}\\
&\hspace{0.35cm}0-2x^3+4x^2\\
&\hspace{0.10cm}-\underline{(-2x^3+2x^2)}\\
&\hspace{1.5cm}0+2x^2-3x\\
&\hspace{1.6cm}-\underline{(2x^2-2x)}\\
&\hspace{2.7cm}0-\;\;x+1\\
&\hspace{2.56cm}-\underline{(-\;x+1)}\\
&\hspace{4.4cm}0\\
\end{align*}
Dann gilt:\\
\begin{align*}
\frac{P_2(x)}{(x-1)}=x^3-2x^2+2x-1
\end{align*}
\begin{align*}
\Leftrightarrow P_2(x)=(x^3-2x^2+2x-1)(x-1)\\
\end{align*}
\paragraph{Aufgabe zur Nullstellenberechnung Teil III}$  $\vspace*{3mm}\\
\href{https://youtu.be/I40ZrXwe9rc}{Video 10}\\
Faktorisiere $ P_3(x)=x^3-2x^2+2x-1 $\\
Suche Nullstelle:\\
\begin{align*}
P_3(1)=1-2+2-1=0\\
\end{align*}
Also ist $ (x-1) $ teiler von $ P_2(x) $\\
\begin{align*}
&(x^3-2x^2+2x-1)\div (x-1)=x^2-x+1\\
-&\underline{(x^3-\;\;x^2)}\\
&\hspace{0.35cm}0-\;\;x^2+2x\\
&\hspace{0.3cm}-\underline{(-x^2+\;\;x)}\\
&\hspace{1.55cm}0+\;\;x-1\\
&\hspace{1.9cm}-\underline{(x-1)}\\
&\hspace{3.4cm}0\\
\end{align*}
Dann gilt:\\
\begin{align*}
\frac{P_3(x)}{(x-1)}=x^2-x+1
\end{align*}
\begin{align*}
\Leftrightarrow P_3(x)=(x^2-x+1)(x-1)\\
\end{align*}
Faktorisiere zuletzt: $ P_4(x)=x^2-x+1 $\\
Mit quatratischer Ergänzung oder p-q-Formel:\\
\begin{align*}
x^2-x+1&=0\\
x^2-x+\frac{1}{4}-\frac{1}{4}+1&=0\\
\Leftrightarrow\bigg(x-\frac{1}{2}\bigg)^2&=-1+\frac{1}{4}\\
\Leftrightarrow x-\frac{1}{2}&=\pm\sqrt{-\frac{3}{4}}=\sqrt{-1}\pm\sqrt{\frac{3}{4}}\;\;\notin \mathbb{R}\\
\Leftrightarrow x_{6,7}&=\frac{1}{2}\pm \mathrm{i} \sqrt{\frac{3}{4}}\\
x_6&=\frac{1+\mathrm{i}\sqrt{3}}{2}\\
x_7&=\frac{1-\mathrm{i}\sqrt{3}}{2}\\
\end{align*}
Zusammenfassung der Faktorisierung von $ P(x) $\\
Faktorisierung von $ P(x) $ mit komplexen Argument:\\
\begin{align*}
p(x)&=x^7-4x^6+8x^5-11x^4+11x^3-8x^2+4x-1\\
&=(x^4-3x^3+4x^2-3x+1)(x-1)(x-\mathrm{i})(x+\mathrm{i})\\
&=(x^3-2x^2+2x-1)(x-1)(x-1)(x-\mathrm{i})(x+\mathrm{i})\\
&=(x^2-x+1)(x-1)(x-1)(x-1)(x-i)(x+i)\\
&=\bigg(x-\bigg(\frac{1+\mathrm{i}\sqrt{3}}{2}\bigg)\bigg)\bigg(x-\bigg(\frac{1-\mathrm{i}\sqrt{3}}{2}\bigg)\bigg)(x-1)^3(x-\mathrm{i})(x+\mathrm{i})\\
\end{align*}
Faktorisierung von $ P(x) $ mit reellen Argument:\\
\begin{align*}
p(x)&=x^7-4x^6+8x^5-11x^4+11x^3-8x^2+4x-1\\
&=(x^2-x+1)(x^2+1)(x-1)^3\\
\end{align*}
\subsection{Koeffizientenvergleich}
\href{https://youtu.be/Zipnj1HU4aU}{Video 11}\\
Seien $ p(x)=a_0+a_1x+\ldots+a_nx^n$ und $q(x) =b_0+b_1x+\ldots +b_kx^k $ sowie $ k\leq n $.\\
Stimmen $ p $ und $ q $ auf $ n $ Stellen überein, so gilt $ n=k $ und $ a_i=b_i $ für $ i=0, \cdots ,n $.\vspace*{2mm}\\
Wenn gilt $ p(x_i)=q(x_i) $ für $ i\in\{0,\cdots ,n\} $, dann gilt $ p(x)=q(x) $\\
Die $ x_i $ und $ x_j $ sollen unterschiedlich sein, weil man sonst immer den gleichen Wert einnehmen kann und dies bringt keine Aussage bis $ n $. Deswegen $ x_i\neq x_j $\vspace*{2mm}\\
Beispiel:\\
Schreibe das Polynom $ 3x^2+12 $ als Summe der Polynome $ x^2-3x $ und $ x^2+4 $.\\
Bestimme:\\
\begin{align*}
&A,B \in \mathbb{R}\;\;\text{ mit}\\
&3x^2+12=A(x^2-3x)+B(x^2+4)\\
\end{align*}
Dies ist:\\
\begin{align*}
\Leftrightarrow 3x^2+0\cdot x+12&=Ax^2-3Ax+Bx^2+4B\\
&=(A+B)x^2-3Ax+4B\\
\end{align*}
Es muß übereinstimmen:\\
\begin{align*}
\textcolor{red}{3}x^2\textcolor{blue}{+0}\cdot x \textcolor{green}{+12}=\textcolor{red}{(A+B)}x^2\textcolor{blue}{-3A}x+\textcolor{green}{4B}\\
\end{align*}
Man bekommt dann 3 Relationen und bestimmt dann $ A $ und $ B $.\\
Dies führt auf 3 LGS (lineares Gleichungssystem):\\
\begin{align*}
3&=A+B&\Rightarrow B=3\\
0&=-3A&\Rightarrow A=0\\
12&=4B&\Rightarrow 12=12\\
\end{align*}
Also gilt:\\
\begin{align*}
3x^2+12=0(x^2-3x)+3(x^2+4)\\
\end{align*}
\subsubsection{Aufgabe zum Koeffizientenvergleich}
\href{https://youtu.be/lLfxDN2AW4k}{Video 12}\\
Bestimme zu folgenden Eigenschaften das entsprechende Polynom $ p(x) $.\\
\begin{quote}
	(i)\;\; $ p(x)=p(-x)$ für alle $ x\in \mathbb{C} $\\
	(ii)\;\;$ p(2\mathrm{i})=p(7)=0 $ und $ p(0)=2 $\\
	(iii) $ \deg(p)\leq 4 $\\
\end{quote}
Bestimmung von $ p(x) $:\\
Wegen:\\
\begin{align*}
p(x)=p(-x)
\end{align*}
gilt:\\
\begin{align*}
p(2i)=p(-2i)=p(7)=p(-7)=0
\end{align*}
d.h.: Die Nullstellen von $ p(x) $ sind:\\
\begin{align*}
&x_1=2i\\
&x_2=-2i\\
&x_3=7\\
&x_4=-7\\
\end{align*}
Da $ \deg(p)\leq 4 $ hat $ p $ max. 4 Nullstellen.\\
Dann git wegen den Fundamentalsatz der Algebra:\\
\begin{align*}
p(x)&=\mathrm{k}(x-2\mathrm{i})(x+2\mathrm{i})(x+7)(x-7)\\
&=\mathrm{k}(x^2+4)(x^2-49)=\mathrm{k}(x^4-49x^2+4x^2-196)\\
&=\mathrm{k}(x^4-45x^2-196)\\
&=\mathrm{k}x^4-45\mathrm{k}x^2-\mathrm{k}196\\
\end{align*}
Bestimmung von k:\\
Es gilt:\\
\begin{align*}
&p(0)=2\Rightarrow \mathrm{k}\cdot 0^4 -45\mathrm{k}\cdot 0^2-\mathrm{k}196=2\\
&\Leftrightarrow -\mathrm{k}196=2\\
&\Leftrightarrow \mathrm{k}= -\frac{1}{98}\\
\end{align*}
Fazit:\\
\begin{align*}
p(x)=-\frac{1}{98}x^4+\frac{45}{98}x^2+2\\
\end{align*}
\subsection{Partialbruchzerlegung rationaler Funktionen}
\href{https://youtu.be/imQZqfJI8yI}{Video 13}\\
Eine rationale Funktion $ f(x)=\frac{p(x)}{q(x)} $ mit $ x\in \mathbb{C} $, wobei $ p $ und $ q $ Polynome sind und  $ \deg(q)>\deg(p) $  ist.\\
Wir wollen $ f(x) $ als Summe einfacher Summanden der Form:\\
\begin{align*}
f(x)=\frac{ax+b}{(x-x_i)^n}
\end{align*}
mit $ x_i $ sind Nullstellen von $ q(x) $ schreiben.\\
Das Ziel ist eine rationale Funktion in eine einfache rationale Funktion zu schreiben.\vspace*{2mm}\\
Beispiel:\\
\begin{align*}
f(x)=\frac{5x^2-4x+7}{x^3+x^2-5x+3}=\frac{5x^2-4x+7}{(x-1)^2(x+3)}=\frac{1}{(x-1)}+\frac{2}{(x-1)^2}+\frac{4}{(x+3)}\\
\end{align*}
Es gibt drei Fälle:\\
\underline{1.Fall:}\\
\begin{quote}
	$ q(x) $ zerfällt vollständig $ q(x)=(x-x_1)\cdot\ldots\cdot(x-x_n) $ und $ x_i $ verschieden.\\
\end{quote}
\underline{2.Fall:}\\
\begin{quote}
	$ q(x) $ zerfällt vollständig $ q(x)=(x-x_1)\cdot\ldots\cdot(x-x_n) $ aber $ x_i $ nicht verschieden.\\
\end{quote}
\underline{3.Fall:}\\
\begin{quote}
	$ q(x) $ zerfällt nicht vollständig. \\
\end{quote}
\subsubsection{PBZ: Nenner zerfällt vollständig mit verschiedenen Nullstellen.}
\href{https://youtu.be/MYsd2kAmsk0}{Video 14}\\
$ f(x)=\frac{p(x)}{q(x)} $, $ q(x) $ zerfällt vollständig in verschiedenen Linearfaktoren.\\
Beispiel:\\
Bestimme die PBZ (Partialbruchzerlegung) für:\\
\begin{align*}
f(x)=\frac{7x}{x^3-x^2+x-1}
\end{align*}
Nullstellen von $ x^3-x^2+x-1=0 $ sind:\\
\begin{align*}
&x_1=1\\
&x_2=\mathrm{i}\\
&x_3=-\mathrm{i}\\
\end{align*}
Es muß gelten:\\
\begin{align*}
f(x)=\frac{7x}{x^3-x^2+x-1}=\frac{7x}{(x-1)(x-\mathrm{i})(x+\mathrm{i})}=\frac{A}{(x-1)}+\frac{B}{(x-\mathrm{i})}+\frac{C}{(x+\mathrm{i})}\\
\end{align*}
Bestimme A,B und C:\\
Vorgehensweise:\\
Erklärung der Zuhaltemethode:\\
Mulitipliziere die Gleichung mit $ (x-1) $:\\
\begin{align*}
\frac{7x}{(x-1)(x-\mathrm{i})(x+\mathrm{i})}&=\frac{A}{(x-1)}+\frac{B}{(x-\mathrm{i})}+\frac{C}{(x+\mathrm{i})}\;\;|\cdot (x-1)\\
\frac{7x(x-1)}{(x-1)(x-\mathrm{i})(x+\mathrm{i})}&=\frac{A(x-1)}{(x-1)}+\frac{B(x-1)}{(x-\mathrm{i})}+\frac{C(x-1)}{(x+\mathrm{i})}\\
\frac{7x}{(x-\mathrm{i})(x+\mathrm{i})}&=A+\frac{B(x-1)}{(x-\mathrm{i})}+\frac{C(x-1)}{(x+\mathrm{i})}\\
\end{align*}
Setze die Nullstelle $ x_1=1 $ ein.\\
\begin{align*}
\frac{7\cdot 1}{(1-\mathrm{i})(1+\mathrm{i})}&=A+\frac{B(1-1)}{(1-\mathrm{i})}+\frac{C(1-1)}{(1+\mathrm{i})}\\
\frac{7}{(1-\mathrm{i})(1+\mathrm{i})}&=A\\
\end{align*}
Berechnung von A,B und C mit der Zuhaltemethode.\\
\begin{align*}
\frac{7x}{\textcolor{red}{(x-1)}\textcolor{blue}{(x-\mathrm{i})}\textcolor{green}{(x+\mathrm{i})}}&=\frac{A}{(x-1)}+\frac{B}{(x-\mathrm{i})}+\frac{C}{(x+\mathrm{i})}\;\;|\cdot (x-1)\\
\end{align*}
Durch Zuhalten von $ \textcolor{red}{(x-1)} $ und für $ x=1 $ einsetzen ergibt dann A:\\
\begin{align*}
\frac{7\cdot1}{(1-\mathrm{i})(1+\mathrm{i})}=\frac{7}{(1-\mathrm{i})^2}=\frac{7}{1+1}=\frac{7}{2}=A\\
\end{align*}
Durch Zuhalten von $ \textcolor{blue}{(x-\mathrm{i})} $ und für $ x=\mathrm{i} $ einsetzen ergibt dann B:\\
\begin{align*}
\frac{7\cdot \mathrm{i}}{(\mathrm{i}-1)(\mathrm{i}+1)}=\frac{7\mathrm{i}}{(\mathrm{i}-1)2\mathrm{i}}=\frac{7\mathrm{i}}{-2-2\mathrm{i}}=B\\
\end{align*}
Durch Zuhalten von $ \textcolor{green}{(x+\mathrm{i})} $ und für $ x=-\mathrm{i} $ einsetzen ergibt dann C:\\
\begin{align*}
\frac{7\cdot(-\mathrm{i})}{(-\mathrm{i}-1)(-\mathrm{i}-\mathrm{i})}=\frac{-7\mathrm{i}}{(-\mathrm{i}-1)(-2\mathrm{i})}=\frac{-7\mathrm{i}}{-2+2\mathrm{i}}=C\\
\end{align*}
Fazit:\\
\begin{align*}
\frac{7x}{x^3-x^2+x-1}=\frac{\frac{7}{2}}{(x-1)}+\frac{\frac{7\mathrm{i}}{-2-2\mathrm{i}}}{(x-\mathrm{i})}+\frac{\frac{-7\mathrm{i}}{-2+2\mathrm{i}}}{(x+\mathrm{i})}\\
\end{align*}
\subsubsection{PBZ: Nenner hat mehrfache Nullstellen}
\href{https://youtu.be/4sYdi6gqWZo}{Video 15}\\
$ f(x)=\frac{p(x)}{q(x)} $, $ q(x) $ zerfällt vollständig in Linearfaktoren mit mehrfachen Nullstellen.\\
Beispiel:\\
Bestimme die PBZ für:\\
\begin{align*}
f(x)=\frac{1}{x^3(x+1)^2}\\
\end{align*}
Die Nullstellen von $ x^3(x+1)^2 =0$ sind:\\
\begin{align*}
x_1, x_2, x_3& = 0\\
x_4, x_5& = -1\\
\end{align*}
Es muß gelten:\\
\begin{align*}
\frac{1}{x^3(x+1)^2}=\frac{A_1}{x}+\frac{A_2}{x^2}+\frac{A_3}{x^3}+\frac{B_1}{(x+1)}+\frac{B_2}{(x+1)^2}\\
\end{align*}
Bestimmung von $ A_1, A_2, A_3, B_1$ und $B_2$:\\
Für die höchsten Potenzen ($ A_3, B_2 $) funktioniert die Zuhaltemethode:\\
\begin{align*}
A_3=\frac{1}{(0+1)^2}=1\\
\end{align*}
\begin{align*}
B_2=\frac{1}{(-1)^3}=-1\\
\end{align*}
Bleibt noch $ A_1, A_2 $ und $ B_1 $ zu bestimmen:\\
Es gilt nun:\\
\begin{align*}
\frac{1}{x^3(x+1)^2}=\frac{A_1}{x}+\frac{A_2}{x^2}+\frac{1}{x^3}+\frac{B_1}{(x+1)}-\frac{1}{(x+1)^2}\\
\end{align*}
Die Gleichung mit dem Nenner: $ x^3(x+1)^2 $ mulitiplizieren.\\
\begin{align*}
1&=\frac{A_1x^3(x+1)^2}{x}+\frac{A_2x^3(x+1)^2}{x^2}+\frac{x^3(x+1)^2}{x^3}+\frac{B_1x^3(x+1)^2}{(x+1)}-\frac{x^3(x+1)^2}{(x+1)^2}\\
&=A_1x^2(x+1)^2 +A_2x(x+1)^2+(x+1)^2+B_1x^3(x+1)-x^3\\
\end{align*}
Es gilt:\\
\begin{align*}
0x^4+0x^3+0x^2+0x+1&=A_1x^2(x+1)^2 +A_2x(x+1)^2+(x+1)^2+B_1x^3(x+1)-x^3\\
&=A_1x^2(x^2+2x+1)+A_2x(x^2+2x+1)+x^2+2x+1+B_1x^3(x+1)-x^3\\
&=A_1(x^4+2x^3+x^2)+A_2(x^3+2x^2+x)+x^2+2x+1+B_1(x^4+x^3)-x^3\\
\end{align*}
Jetzt die Gleichung nach $ x $ Potenzen sortieren:\\
\begin{align*}
0x^4+0x^3+0x^2+0x+1&=(A_1+B_1)x^4+(2A_1+A_2+B_1-1)x^3+(A_1+2A_2+1)x^2+1\\
\end{align*}
%\subsubsection{ PBZ: Nenner hat mehrfache Nullstellen Teil II}
\href{https://youtu.be/_I2IcH8-JMw}{Video 16}\\
Mit Koeffizientvergleich bekommt man folgendes LGS:\\
\begin{align}
0&=A_1+B_1 \label{a1}\\
0&=2A_1+A_2+B_1-1 \label{a2}\\
0&=A_1+2A_2+1 \label{a3}\\
0&=A_2+2\label{a4}\\
1&=1
\end{align}
Mit \ref{a4} :\\
\begin{align*}
&0=A_2+2\\
&\Leftrightarrow A_2=-2\\
\end{align*}
$ A_2=-2 $ in \ref{a3} einsetzen:\\
\begin{align*}
&0=A_1+2A_2+1=A_1+2(-2)+1\\
&\Leftrightarrow A_1=3\\
\end{align*}
$ A_1=3 $ in\ref{a1} einsetzen:\\
\begin{align*}
&0=A_1+B_1=3+B_1\\
&\Leftrightarrow B_1=-3\\
\end{align*}
Fazit:\\
\begin{align*}
\frac{1}{x^3(x+1)^2}=\frac{3}{x}-\frac{2}{x^2}+\frac{1}{x^3}-\frac{3}{(x+1)}-\frac{1}{(x+1)^2}\\
\end{align*}

\subsubsection{PBZ: Nenner zerfällt nicht vollständig }
\href{https://youtu.be/EepB1ozV-fQ}{Video 17}\\
$ f(x)=\frac{p(x)}{q(x)} $, $ q(x) $ zerfällt nicht vollständig in Linearfaktoren.\\
Kann wegen des Fundamentalsatzes der Algebra nur auftreten, wenn nur reelle Nullstellen für $ q(x) $ zugelassen sind.\\
Beispiel:\\
Bestimme die PBZ für:\\
\begin{align*}
f(x)=\frac{2x-1}{(x-2)(x^2+2)}\;\;\text{mit }\;\;x\in\mathbb{R}\\
\end{align*}
Die Nullstellen von $ (x-2)(x^2+2)=0 $ sind:\\
\begin{align*}
x_1=2
\end{align*}
Bemerkung:\\
Die Nullstellen von $ x^2+2=0 $ sind:\\
\begin{align*}
x_2=\mathrm{i}\sqrt{2}\\
x_3=-\mathrm{i}\sqrt{2}\\
\end{align*}
D.h.: $ x^2+2 $ zerfällt nicht weiter, da $ x_2, x_3 $ komplexe Zahlen sind ($x_1,x_2 \notin\mathbb{R} $). \\
Wenn $ x\in \mathbb{C} $, dann zerfällt das Polynom komplett.\vspace*{2mm}\\
Es muß gelten:\\
\begin{align*}
f(x)=\frac{2x-1}{(x-2)(x^2+2)}=\frac{A}{(x-2)}+\frac{Bx+C}{(x^2+2)}\\
\end{align*}
Bemerkung:\\
$ Bx+c $ weil $ x^2+2 $ nicht weiter zerfällt.\\
%\subsubsection{PBZ: Nenner zerfällt nicht vollständig Teil II }
\href{https://youtu.be/FvWtzBpsLvk}{Video 18}\\
Berechnung von $ A,B $ und $ C $.\\
Bestimmung von $ A $ mit der Zuhaltemethode:\\
\begin{align*}
A=\frac{2\cdot 2-1}{2^2+2}=\frac{3}{6}=\frac{1}{2}\\
\end{align*}
Bestimmung von $ B $ und $ C $:\\
Es gilt nun:\\
\begin{align*}
\frac{2x-1}{(x-2)(x^2+2)}=\frac{\frac{1}{2}}{(x-2)}+\frac{Bx+C}{(x^2+2)}\\
\end{align*}
Die Gleichung mit dem Nenner $ (x-2)(x^2+2) $ multiplizieren.\\
\begin{align*}
\frac{(2x-1)(x-2)(x^2+2)}{(x-2)(x^2+2)}&=\frac{\frac{1}{2}(x-2)(x^2+2)}{(x-2)}+\frac{(Bx+C)(x-2)(x^2+2)}{(x^2+2)}\\
2x-1&=\frac{1}{2}(x^2+2)+(Bx+C)(x-2)\\
&=\frac{1}{2}x^2 +1+Bx^2-2Bx+Cx-2C\\
0x^2+2x-1&=\bigg(\frac{1}{2}+B\bigg)x^2+(C-2B)x+1-2C\\
\end{align*}
Mit Koeffizientvergleich bekommt man folgendes LGS:\\
\begin{align*}
&0=\frac{1}{2}+B\\
&\Leftrightarrow B=-\frac{1}{2}\\
&-1=1-2C\\
&\Leftrightarrow C=1\\
\end{align*}
Fazit:\\
\begin{align*}
&A=\frac{1}{2}\\
&B=-\frac{1}{2}\\
&C=1\\
\end{align*}
Daraus folgt:\\
\begin{align*}
f(x)=\frac{2x-1}{(x-2)(x^2+2)}=\frac{\frac{1}{2}}{(x-2)}+\frac{(-\frac{1}{2})x+1}{(x^2+2)}
\end{align*}
\subsubsection{PBZ: Nennergrad ist kleiner gleich als der Zählergrad }
\href{https://youtu.be/ovgTWyr47GY}{Video 19}\\
\begin{align*}
f(x)\frac{p(x)}{q(x)}\;\;\text{ mit}\;\;\deg(p)\geq \deg (q)\\
\end{align*}
In diesen Fall verwendet wir die Polynomdivision.\\
Wähle ein Polynom $ h(x) $ und $ r(x) $ mit.\\
\begin{align*}
p(x)=h(x)\cdot q(x)+r(x)\;\;\text{mit } \deg(r)<\deg(q)\\
\end{align*}
\newpage
\noindent
Bemerkung:\\
\begin{quote}
	

Das Polynom: $ \frac{7}{2} $\\
Addition und Multiplikation bei Polynomen gilt genau so wie bei Zahlen.\\
Wenn man den Bruch $ \frac{7}{2} $ teilt, dann ist der Rest immer kleiner als der Nenner.\\
\begin{align*}
7=3\cdot 2+1
\end{align*}

 Wähle $h(x)=3  $ und $ p(x)=7 $ dann ergibt sich $ r(x)=1 $ mit der Bedingung $\deg(r)<\deg(q)  $.\\
\begin{align*}
p(x)&=h(x)q(x)+r(x)\\
7&=\;\;3\;\;\cdot\; 2 \;\;+\;1\\
\end{align*}
Rest=1 ist imme kleiner als der Nenner=2.
\end{quote}
Nun teilt man die das Polynom $ p(x)=h(x)q(x)+r(x) $ mit $ q(x) $.\\
\begin{align*}
f(x)=\frac{p(x)}{q(x)}=h(x)+\frac{r(x)}{q(x)}\\
\end{align*}
Nun, da $ \deg(r)<\deg(q) $ gilt können wir die PBZ auf $ \frac{r(x)}{q(x)}$ anwenden.\vspace*{2mm}\\
Beispiel:\\
\begin{align*}
f(x)=\frac{p(x)}{q(x)}=\frac{x^2+x+1}{x-2}\\
\end{align*}
Wegen $ \deg(p)=2>1=\deg(q) $ muß  zunächst die Polynomdivision durchgeführt werden.\\
\begin{align*}
&(x^2+x+1)\div(x-2)=x+3\\
-&\underline{(x^2-2x)}\\
&\hspace{1.1cm} 3x+1\\
&\hspace{0.5cm}-\underline{(3x-6)}\\
&\hspace{2.2cm}7\\
\end{align*}
\begin{align*}
&h(x)=x+3\\
&r(x)=7\\
\end{align*}
Fazit:\\
\begin{align*}
f(x)=\frac{p(x)}{q(x)}=\frac{x^2+x+1}{x-2}=x+3+\frac{7}{x-2}\\
\end{align*}
\subsubsection{Partialbruchzerlegung: Anleitung }
\href{https://youtu.be/ADKDURyGOmg}{Video 20}\\
Anleitung für eine Partialbruchzerlegung von $ f(x)=\frac{p(x)}{q(x)} $\\
1. Bestimmung von $ \deg(p) $ und $ \deg(q) $\vspace*{2mm}\\
2.1 $ \deg(p)< \deg(q) $\\
\begin{quote}
	$\blacktriangleright$ Faktorisierung von $ q(x) $\\
	\begin{quote}
		a. $ q(x) $ zerfällt vollständig in Linearfaktoren\\
		\begin{quote}
			\textbullet \;\;mit einfacher Nullstelle $ \Rightarrow $ Zuhaltemethode\\
			\textbullet \;\;mit mehrfacher Nullstelle $ \Rightarrow $ mit Koeffizientenvergleich $ \Rightarrow $ LGS\\
		\end{quote}
	    b. $ q(x) $ zerfällt nicht vollständig\\
        \begin{quote}
	    Betrachte Ansatz mit $ \frac{Bx+C}{x-x_i} \Rightarrow$ Koefizientenvergleich $ \Rightarrow $ LGS\\
	    \end{quote}
	\end{quote}
\end{quote}
2.2 $ \deg(p)> \deg(q) $\\
\begin{quote}
	$\blacktriangleright$ Bestimme mit Polynomdivision
	 \begin{align*}
	  &h(x) \;\;\text{und }\;\;r(x)\\
	  &\frac{p(x)}{q(x)}=h(x)+\frac{r(x)}{q(x)}\;\;\text{mit }  \deg(r)<\deg(q) \\
	  &\text{dann weiter mit 2.1 }
	 \end{align*}
	
\end{quote}